%%%%%%%%%%%%%%%%%%%%%%%%%%%%%%%%%%%%%%%%%%%%%%%%%%%%%%%%%%
% FILE: chapter10_buenas_practicas.tex
%
% PURPOSE
% -------------------------------------------------------
% Provides best practices for using and maintaining
% the template in academic and technical projects.
%
%
% ARCHITECTURE ROLE
% -------------------------------------------------------
% Final chapter of the template documentation.
%
% Describes recommended practices for structuring
% documents and maintaining the template.
%
%
% DEPENDENCIES
% -------------------------------------------------------
% Entire template structure
%
%
% NOTES FOR DEVELOPERS
% -------------------------------------------------------
% This chapter should remain conceptual.
%
% It documents guidelines rather than specific
% implementation details.
%
%%%%%%%%%%%%%%%%%%%%%%%%%%%%%%%%%%%%%%%%%%%%%%%%%%%%%%%%%%


\chapter{Buenas prácticas}

Este capítulo recoge recomendaciones para utilizar
la plantilla de forma eficiente y mantener la
consistencia del documento.



% --------------------------------------------------
% SEPARACIÓN DE RESPONSABILIDADES
% --------------------------------------------------

\section{Separación de responsabilidades}

\begin{definitionbox}[Separación de responsabilidades]
La separación de responsabilidades consiste en dividir
un sistema en componentes independientes, cada uno
encargado de una función específica.
\end{definitionbox}

La plantilla está organizada siguiendo este principio.

Cada carpeta tiene un propósito concreto.



\begin{lstlisting}[caption={Arquitectura de carpetas},label={code:structure}]
00_config/
01_frontmatter/
02_chapters/
03_figures/
04_bibliography/
05_appendices/
06_styles/
07_acronyms/
\end{lstlisting}

\source{Plantilla}



\begin{notebox}
Mantener esta separación facilita el mantenimiento
y la reutilización de la plantilla.
\end{notebox}



% --------------------------------------------------
% QUÉ ARCHIVOS MODIFICAR
% --------------------------------------------------

\section{Archivos que el usuario debe modificar}

Durante la redacción del documento, el usuario
debe trabajar principalmente en los siguientes
archivos y carpetas.

\begin{itemize}

\item \texttt{00\_config/config.tex}
\item \texttt{02\_chapters/}
\item \texttt{03\_figures/}
\item \texttt{04\_bibliography/}
\item \texttt{07\_acronyms/}

\end{itemize}



% --------------------------------------------------
% QUÉ ARCHIVOS NO MODIFICAR
% --------------------------------------------------

\section{Archivos que no deberían modificarse}

Los archivos de estilo y configuración avanzada
no deberían modificarse salvo que se comprenda
su funcionamiento.

\begin{itemize}

\item \texttt{06\_styles/}
\item \texttt{packages.tex}
\item \texttt{main.tex}

\end{itemize}



\begin{warningbox}
Modificar archivos de estilo sin comprender su
función puede provocar errores de compilación
o inconsistencias tipográficas.
\end{warningbox}



% --------------------------------------------------
% ORGANIZACIÓN DE CAPÍTULOS
% --------------------------------------------------

\section{Organización de capítulos}

Cada capítulo debe almacenarse en un archivo
independiente dentro de la carpeta:

\texttt{02\_chapters}



Ejemplo:

\begin{lstlisting}[caption={Organización de capítulos},label={code:chapters}]
chapter01_introduccion.tex
chapter02_arquitectura.tex
chapter03_metodologia.tex
chapter04_resultados.tex
\end{lstlisting}

\source{Plantilla}



\begin{examplebox}
Esta organización permite trabajar en documentos
grandes de forma modular y facilita la colaboración
entre varios autores.
\end{examplebox}



% --------------------------------------------------
% NOMBRES DE ARCHIVOS
% --------------------------------------------------

\section{Convenciones de nombres}

Se recomienda utilizar nombres de archivo
descriptivos y consistentes.

Ejemplo:

\begin{lstlisting}[caption={Convención de nombres},label={code:naming}]
chapter01_introduccion.tex
chapter02_estado_del_arte.tex
chapter03_metodologia.tex
\end{lstlisting}

\source{Plantilla}



Esto facilita identificar rápidamente el
contenido de cada archivo.



% --------------------------------------------------
% ERRORES COMUNES
% --------------------------------------------------

\section{Errores comunes}

Durante el uso de la plantilla pueden aparecer
algunos errores frecuentes.



\begin{warningbox}
Olvidar ejecutar \texttt{biber} provocará que
las referencias bibliográficas no aparezcan
en el documento final.
\end{warningbox}



\begin{warningbox}
Insertar figuras o tablas sin utilizar las macros
definidas por la plantilla puede generar
inconsistencias de formato.
\end{warningbox}



\begin{warningbox}
Editar directamente \texttt{main.tex} para escribir
contenido rompe la arquitectura modular del proyecto.
\end{warningbox}



% --------------------------------------------------
% MANTENIMIENTO
% --------------------------------------------------

\section{Mantenimiento de la plantilla}

Para mantener la plantilla organizada se recomienda:

\begin{itemize}

\item documentar los cambios realizados
\item mantener comentarios claros en los archivos
\item respetar la estructura de carpetas
\item evitar modificaciones innecesarias en los estilos

\end{itemize}



\begin{examplebox}
Una plantilla bien mantenida puede reutilizarse
en múltiples proyectos académicos y técnicos.
\end{examplebox}
